\section{Capsule two: keeping infinitely many indices}\label{ch17:sec:pigeonhole}
A sequence $x_0,x_1,x_2,\ldots$ has one term for each index $n\in\NN$, so it has infinitely many indices, even when its values repeat. The sequence $x_n=(-1)^n$ takes only the two values $1$ and $-1$, yet the value $1$ occurs at the infinitely many indices $0,2,4,\ldots$. A constant sequence takes the same value at every index. For this reason the proof keeps track of sets of \emph{indices}, not sets of points. When we say that a ball contains infinitely many terms of a sequence, we mean that it contains $x_n$ for infinitely many indices $n$, whether or not these terms are different points.

Recall from Section~\ref{ch03:sec:sets} that a set is finite if its members can be written down in a list that ends. A set that is not finite is called \emph{infinite}. For example, $\NN$ is infinite, and so is the set $\set{0,2,4,\ldots}$ of even natural numbers.

\par\noindent\begin{minipage}{\linewidth}
\begin{lemma}[Infinite pigeonhole principle]\label{ch17:lem:pigeonhole}
Let $S$ be an infinite set, and let $A_1,\ldots,A_k$ be finitely many sets whose union contains $S$. Then at least one of the sets $S\cap A_1,\ldots,S\cap A_k$ is infinite.
\end{lemma}
\end{minipage}\par

\begin{proof}
Every member of $S$ lies in at least one of the sets $A_i$, so
\[
 S=(S\cap A_1)\cup(S\cap A_2)\cup\cdots\cup(S\cap A_k).
\]
Suppose, for contradiction, that each of the sets $S\cap A_i$ is finite. Write down the members of $S\cap A_1$, then the members of $S\cap A_2$, and so on up to $S\cap A_k$. The result is a list that ends, and every member of $S$ appears in it at least once. A member that lies in several of the sets $A_i$ appears several times; crossing out the repetitions leaves a list of the members of $S$ that still ends. So $S$ is finite, which contradicts the assumption that $S$ is infinite.
\end{proof}

Compare this with the pigeonhole principle of Proposition~\ref{ch07:prop:pigeonhole}, which says that when many objects are placed in a few boxes, some box must hold many of them. Here the objects are the members of $S$ and the boxes are the sets $A_1,\ldots,A_k$: infinitely many objects in finitely many boxes force some box to hold infinitely many. The boxes are allowed to overlap. The proof only needed to show that $S$ would be finite, and an object that lies in two boxes merely appears twice in the list.

For example, take $x_n=(-1)^n$ and the two balls $B(1,1/2)$ and $B(-1,1/2)$ on the real line. Let $A_1$ be the set of indices $n$ with $x_n\in B(1,1/2)$, which is the set of even numbers, and let $A_2$ be the set of indices with $x_n\in B(-1,1/2)$, which is the set of odd numbers. Their union is $\NN$, and the lemma, applied with $S=\NN$, promises that at least one of them is infinite. Here both are. The lemma guarantees at least one, but it does not say which.

The proof of the theorem uses the lemma at every stage in exactly this way. The set $S$ is the current infinite set of indices. The space is covered by finitely many balls, and $A_i$ is the set of indices $n$ for which $x_n$ lies in the $i$th ball. Since the balls cover the space, every term $x_n$ lies in at least one of them, so the union of the sets $A_i$ contains every index. The lemma then provides one ball that contains $x_n$ for infinitely many indices $n$ in $S$.
